\section{Public key encryption schemes}\label{sec:prelims-encryption}

Public key encryption schemes enable confidential communication by allowing
anyone possessing a public key to encrypt messages that can only be decrypted
using the corresponding secret key.

\paragraph{Syntax.} A public key encryption scheme consists of the following
algorithms:

\begin{description}

  \item[$(\sk, \pk) \sample \PKEgen(\secparam)$ : ] The key generation algorithm
outputs secret decryption key $\sk$ and public encryption key $\pk$.

  \item[$\ciph \sample \PKEencrypt(\pk, \msg)$ : ] The encryption algorithm
encrypts message $\msg \in \msgspace$ using public key $\pk$, giving output
ciphertext
$\ciph$.

  \item[$\msg \leftarrow \PKEdecrypt(\sk, \ciph)$ : ] The decryption algorithm
decrypts ciphertext $\ciph$ using secret key $\sk$, giving output message
$\msg$.

\end{description}

Public key encryption schemes satisfy correctness and indistinguishability. We
define correctness and indistinguishability under chosen-plaintext attack
below. We use the standard syntax and security notions for public key
encryption~\cite{KatLin14}.

\paragraph{Correctness.} A public key encryption scheme is correct if, for every
key pair $(\sk,\pk)$ generated by $\PKEgen(\secparam)$ and every message $\msg$,
the following holds:
%
\begin{equation*}
%
  \PKEdecrypt(\sk, \PKEencrypt(\pk, \msg)) = \msg.
%
\end{equation*}
%
\paragraph{IND-CPA security.} We define indistinguishability under
chosen-plaintext attack (IND-CPA) security using \Cref{exp:ind-cpa}.

\begin{experimentbox}{Public key encryption schemes: IND-CPA}{ind-cpa}

  \begin{enumerate}

    \item The challenger generates $(\sk, \pk) \sample \PKEgen(\secparam)$ and
sends $\pk$ to $\adv$.

    \item The adversary outputs two equal length messages $(\msg_0,\msg_1)$.

    \item The challenger samples $b \sample \{0,1\}$, computes the challenge
ciphertext $\ciph^\star \sample \PKEencrypt(\pk,\msg_b)$, and sends
$\ciph^\star$ to the adversary.

    \item The adversary outputs a bit $b'$.

  \end{enumerate}

  $\adv$ wins the experiment if $b' = b$.

\end{experimentbox}

The IND-CPA advantage of $\adv$ is defined as follows:
%
\begin{equation*}
%
  \Adv^{\mathsf{ind\text{-}cpa}}_{\PKE, \adv}(\secparam) = \left|
\Pr[\adv\text{ wins}] - \frac{1}{2} \right|.
%
\end{equation*}
%
We say the public key encryption scheme $\PKE$ is IND-CPA secure if this
advantage is negligible for every PPT adversary $\adv$.

Indistinguishability under chosen-ciphertext attack (IND-CCA) is defined by the
same experiment, with the adversary also given an oracle that decrypts any
ciphertext other than $\ciph^\star$. IND-CCA security implies IND-CPA security.
